\documentclass{../template/myrpt}
\ExamName{双光栅测量微弱振动位移}

\newcolumntype{Y}{>{\centering\arraybackslash}X}
\newcolumntype{C}[1]{>{\centering\arraybackslash}m{#1}}
\pgfplotsset{compat=1.18}

\begin{document}
\maketitle

\section{实验目的}
\begin{enumerate}[label=\arabic*. ]
  \item 了解利用光的多普勒频移形成光拍的原理；
  \item 了解精确测量微弱振动位移的一种方法；
  \item 应用双光栅微弱振动测量仪测量音叉振动的微振幅。
\end{enumerate}

\section{实验仪器}

本实验使用的主要仪器为双光栅微弱振动测量仪和示波器。双光栅微弱振动测量仪由半导体激光器、静光栅、与音叉相连的动光栅、光电接收器、驱动与检测电路等部分组成。

\begin{figure}[htbp]
  \centering
  \includegraphics[width=0.72\linewidth]{assets/media/image1.png}
  \caption{双光栅微弱振动测量仪}
\end{figure}

\section{实验原理}

\subsection{移动光学相位光栅的多普勒频移}

当光源和接收器之间存在相对运动时，接收器接收到的光频率会相对于光源频率发生变化，这种现象称为多普勒效应。若把光栅固定在振动音叉上，光栅随音叉作微小周期运动，入射激光经移动光栅衍射后会产生与光栅速度有关的多普勒频移。

设光栅常数为 $d$，光栅相对入射光束的瞬时速度为 $v$。对于同一级衍射光，多普勒频移量可写成
\[
  f_D = \frac{v}{d}.
\]
由此可见，在光栅常数一定时，多普勒频移只与移动光栅的瞬时速度有关，而与光源本身的光频无关。

\subsection{光拍的获得与检测}

光频本身很高，普通电子仪器难以直接测量光频的微小变化。实验中让已发生频移的衍射光与未发生频移或频移不同的光束平行叠加，两束光的频率差会形成光拍。光电接收器具有平方律检波特性，可把光强随时间的变化转换为电信号，示波器上观察到的拍频信号就反映了两束光之间的频率差。

双光栅系统中，一片光栅静止，另一片光栅随音叉振动。激光通过双光栅后，适当选择某一级衍射光进入光电池，就可得到随音叉运动速度变化的光拍信号。调节示波器显示后，可在音叉振动的半个周期 $T/2$ 内计数拍频波形个数。

\subsection{微弱振动位移量的检测}

设音叉带动动光栅沿测量方向作简谐振动
\[
  x = A\sin \omega t,
\]
则动光栅在半个振动周期内从一个极端位置运动到另一个极端位置，位移变化量为 $2A$。拍频信号的瞬时频率与动光栅速度成正比，在半周期内观察到的拍频波形个数 $n$ 满足
\[
  n = \int_0^{T/2} f_D\,\mathrm{d}t
    = \frac{1}{d}\int_0^{T/2} |v|\,\mathrm{d}t
    = \frac{2A}{d}.
\]
因此音叉微振动的振幅为
\[
  A = \frac{nd}{2}.
\]
若光栅密度为 $N$，即 $d=1/N$，则
\[
  A = \frac{n}{2N}.
\]
实验中只要测得半周期内的拍频波形数 $n$，即可计算音叉的微弱振动位移振幅。

\section{实验内容和步骤}

\subsection{仪器连接与光路调节}

\begin{enumerate}[label=\arabic*)]
  \item 打开双光栅微弱振动测量仪和示波器电源，按仪器要求连接检测信号输出端与示波器输入端。
  \item 调节激光器固定架的左右调节和上下调节旋钮，使红色激光依次通过静光栅和动光栅。
  \item 调整光路，使某一级衍射光正好落入光电池前的小孔内，并在示波器上观察到稳定的拍频信号。
\end{enumerate}

\subsection{测量外力驱动音叉的频率--振幅关系}

\begin{enumerate}[label=\arabic*)]
  \item 调节频率粗调旋钮，使外力驱动频率约为 $506\,\si{\hertz}$。
  \item 调节功率旋钮，使驱动电流约为 $40\,\si{\milli\ampere}$。
  \item 细调驱动频率，使音叉在半周期 $T/2$ 内的拍频波形数最多，此时音叉接近谐振状态。
  \item 记录谐振频率 $f_{\text{谐}}$。保持驱动功率不变，在谐振频率附近依次改变驱动频率，记录每个频率下 $T/2$ 内的波形数 $n$。
  \item 根据 $A=nd/2$ 计算音叉振幅，完成表 \ref{table:frequency-amplitude}，并绘制频率--振幅曲线。
\end{enumerate}

\begin{LabRecordTable}
  \centering
  \caption{测量外力驱动音叉的频率--振幅关系（功率 $I=40\,\si{\milli\ampere}$）}
  \label{table:frequency-amplitude}
  \renewcommand{\arraystretch}{1.35}
  \setlength{\tabcolsep}{3pt}
  \small
  \begin{tabularx}{\textwidth}{C{2.2cm} *{7}{Y}}
    \toprule
    频率点 & $f_{\text{谐振}-0.3}$ & $f_{\text{谐振}-0.2}$ & $f_{\text{谐振}-0.1}$ & $f_{\text{谐振}}$ & $f_{\text{谐振}+0.1}$ & $f_{\text{谐振}+0.2}$ & $f_{\text{谐振}+0.3}$ \\
    \midrule
    频率 $f$/\si{\hertz} & 508.700 & 508.800 & 508.900 & 509.000 & 509.100 & 509.200 & 509.300 \\
    波形数 $n$/个 & 4 & 5 & 9 & 30 & 15 & 7 & 5 \\
    振幅 $A$/\si{\milli\metre} & 0.020 & 0.025 & 0.045 & 0.150 & 0.075 & 0.035 & 0.025 \\
    \bottomrule
  \end{tabularx}
\end{LabRecordTable}

\subsection{测量外力驱动音叉的功率--振幅关系}

\begin{enumerate}[label=\arabic*)]
  \item 将驱动频率保持在音叉谐振频率 $f_{\text{谐}}$ 附近。
  \item 依次调节功率旋钮，改变外力驱动功率或驱动电流。
  \item 在每个功率设置下记录半周期 $T/2$ 内的拍频波形数 $n$。
  \item 根据 $A=nd/2$ 计算音叉振幅，完成表 \ref{table:power-amplitude}，绘制功率--振幅曲线，并分析振幅随功率变化的规律。
\end{enumerate}

\begin{LabRecordTable}
  \centering
  \caption{测量外力驱动音叉的功率--振幅关系}
  \label{table:power-amplitude}
  \renewcommand{\arraystretch}{1.35}
  \setlength{\tabcolsep}{4pt}
  \small
  \begin{tabularx}{\textwidth}{C{2.4cm} *{5}{Y}}
    \toprule
    功率 $I$/\si{\milli\ampere} & 5 & 15 & 25 & 35 & 45 \\
    \midrule
    波形数 $n$/个 & 3 & 15 & 22 & 26 & 29 \\
    振幅 $A$/\si{\milli\metre} & 0.015 & 0.075 & 0.110 & 0.130 & 0.145 \\
    \bottomrule
  \end{tabularx}
\end{LabRecordTable}

\section{实验数据处理}

本实验所用光栅密度为 $N=100\,\text{条}/\si{\milli\metre}$，光栅常数为
\[
  d=\frac{1}{N}=0.010\,\si{\milli\metre}.
\]
由微弱振动位移测量公式
\[
  A=\frac{nd}{2}=\frac{n}{2N}
\]
可得本实验中振幅的计算式为
\[
  A=0.005n\,\si{\milli\metre}.
\]
例如，在谐振频率 $f_{\text{谐}}=509.000\,\si{\hertz}$ 时，半周期内拍频波形数为 $n=30$，故
\[
  A=0.005\times 30=0.150\,\si{\milli\metre}.
\]

\subsection{频率--振幅关系}

固定驱动功率 $I=40\,\si{\milli\ampere}$，在谐振频率附近改变驱动频率，得到表 \ref{table:frequency-amplitude} 中的数据。由图 \ref{fig:frequency-amplitude} 可见，音叉振幅在 $509.000\,\si{\hertz}$ 附近达到最大值 $0.150\,\si{\milli\metre}$。当频率低于谐振频率时，振幅随频率升高而增大；当频率高于谐振频率时，振幅随频率继续升高而减小。

\begin{figure}[htbp]
  \centering
  \begin{tikzpicture}
    \begin{axis}[
      width=0.86\linewidth,
      height=5.4cm,
      grid=both,
      xlabel={频率 $f$/\si{\hertz}},
      ylabel={振幅 $A$/\si{\milli\metre}},
      xmin=508.65,
      xmax=509.35,
      ymin=0,
      ymax=0.17,
      xtick={508.7,508.8,508.9,509.0,509.1,509.2,509.3},
      ytick={0,0.025,0.050,0.075,0.100,0.125,0.150},
      yticklabels={0,0.025,0.050,0.075,0.100,0.125,0.150},
      tick label style={font=\small},
      label style={font=\small},
      scaled ticks=false,
    ]
      \addplot[
        mark=*,
        thick,
        blue,
      ] coordinates {
        (508.700,0.020)
        (508.800,0.025)
        (508.900,0.045)
        (509.000,0.150)
        (509.100,0.075)
        (509.200,0.035)
        (509.300,0.025)
      };
    \end{axis}
  \end{tikzpicture}
  \caption{音叉的频率--振幅关系}
  \label{fig:frequency-amplitude}
\end{figure}

\subsection{功率--振幅关系}

保持驱动频率在谐振频率附近，改变驱动电流，得到表 \ref{table:power-amplitude} 中的数据。由图 \ref{fig:power-amplitude} 可见，在本次测量范围内，音叉振幅随驱动电流增大而增大；低电流区间增长较快，较高电流区间增长趋缓。这说明外力驱动增强时输入能量增加，音叉振动更明显，但在接近较大振幅时受阻尼、仪器响应和读数分辨率等因素影响，振幅增量不再严格按相同斜率增加。

\begin{figure}[htbp]
  \centering
  \begin{tikzpicture}
    \begin{axis}[
      width=0.78\linewidth,
      height=5.4cm,
      grid=both,
      xlabel={功率 $I$/\si{\milli\ampere}},
      ylabel={振幅 $A$/\si{\milli\metre}},
      xmin=0,
      xmax=50,
      ymin=0,
      ymax=0.16,
      xtick={5,15,25,35,45},
      ytick={0,0.025,0.050,0.075,0.100,0.125,0.150},
      yticklabels={0,0.025,0.050,0.075,0.100,0.125,0.150},
      tick label style={font=\small},
      label style={font=\small},
      scaled ticks=false,
    ]
      \addplot[
        mark=*,
        thick,
        blue,
      ] coordinates {
        (5,0.015)
        (15,0.075)
        (25,0.110)
        (35,0.130)
        (45,0.145)
      };
    \end{axis}
  \end{tikzpicture}
  \caption{音叉的功率--振幅关系}
  \label{fig:power-amplitude}
\end{figure}

\section{结论和分析}

本实验利用双光栅产生的光拍信号测量音叉的微弱振动位移。实验测得音叉谐振频率为
\[
  f_{\text{谐}}=509.000\,\si{\hertz},
\]
在 $I=40\,\si{\milli\ampere}$ 时谐振点半周期内波形数为 $30$ 个，对应振幅为 $0.150\,\si{\milli\metre}$。频率--振幅曲线表现出明显的谐振特征：驱动频率接近谐振频率时振幅最大，偏离谐振频率后振幅降低。

功率--振幅关系表明，在谐振频率附近，外力驱动电流从 $5\,\si{\milli\ampere}$ 增大到 $45\,\si{\milli\ampere}$ 时，波形数由 $3$ 个增至 $29$ 个，振幅由 $0.015\,\si{\milli\metre}$ 增至 $0.145\,\si{\milli\metre}$。这与外力驱动增强时音叉获得更多能量、振动位移增大的规律一致。

本实验的主要误差来源包括：示波器上拍频波形个数的人工计数误差；谐振频率附近频率调节精度有限；光路未完全调至最佳时会影响光拍信号幅度和清晰度；功率调节和频率读数也受仪器分辨率限制。实验中应使光路稳定、拍频波形清晰，并在谐振点附近细致调节频率，以提高微弱振动位移测量的可靠性。

\insertSignedDataPage

\end{document}
